\section{Objections and replies}
\label{sec:objections}

\subsection{The calculus merely relabels ordinary proof dependencies}

The composition proposition by itself is disciplined bookkeeping, and the
positive spectral lemma is standard.  If the paper claimed a new branch of
mathematics from those facts alone, the objection would succeed.  The added
value lies in the joint constraint imposed by Theorem
\ref{thm:recovery-spectrum} and its countermodels.  A proof ledger that says
only ``lattice estimate implies continuum gap'' does not reveal whether the
failure concerns positivity, dynamical fidelity, physical units, observable
density, or limit order.  The typed framework represents these as independent
obligations, and the lattice cases in Section~\ref{sec:failures} demonstrate
that the distinction selects mathematically different repairs.  The formal
novelty claim is accordingly modest; the methodological novelty claim is
substantive and testable by its effect on proof design.

\subsection{Representations of representations create an infinite regress}

Roberts's hierarchy objection cautions against treating ever-higher structure
as automatically more fundamental \cite{Roberts2011}.  The present calculus
has a stopping criterion: a target physical proposition, its bearer, and a
probe class sufficient for that proposition.  A meta-representation is added
only when an interface lacks a certificate required by the target inference.
Once exact or controlled recovery on a separating domain has been proved, no
further level is demanded by the calculus.  The framework is therefore
goal-relative without being arbitrary; the goal is fixed by the physical
claim, not by a desire for maximal abstraction.

\subsection{Every certificate is itself representation-dependent}

Yes.  Error norms, observable cores, and convergence topologies must be
specified in mathematical language.  The response is not to posit a view
from nowhere but to require compatibility under the next arrow.  A norm used
at one stage is acceptable when the subsequent recovery theorem is continuous
for that norm or when a comparison inequality translates it.  The process can
fail, and that failure is recorded.  Objectivity here means controlled
coherence of claim transport, not representation-free syntax.

\subsection{The framework overconstrains productive physics}

Heuristic analogies, gauge choices, truncations, and uncontrolled limits can
be fruitful discovery tools.  The framework does not forbid them.  It
distinguishes discovery from establishment.  A heuristic arrow may be labeled
as such and used to search for a rigidity engine.  The certificate becomes
mandatory only when the arrow is invoked in a proof of the target physical
claim.  This preserves the flexibility of representation while preventing
its epistemic status from silently changing.

\subsection{Property-relative fidelity makes theory identity fragmentary}

The framework indeed rejects one undifferentiated equivalence verdict for
many scientific purposes.  This does not preclude global equivalence where a
duality or definitional theorem establishes it.  It says that most
constructive transitions are weaker: an RG transformation may be faithful for
long-distance observables and unfaithful for microscopic variables; a gauge
quotient may be exact for invariants and poor for local coordinates.  A global
identity claim is recovered only from a sufficiently rich family of
property-level certificates plus interpretation.  Fragmentation is a feature
of the evidence, not a stipulation of the framework.

\subsection{The ontology is underdetermined}

It is.  The same formal recovery relations can support structural-realist,
entity-realist, empiricist, or pragmatist readings.  This is a limitation only
if the calculus was advertised as a metaphysical deduction.  Its primary
claim concerns warrant and proof methodology.  The conditional ontology in
Section~\ref{sec:ontology} maps possible interpretations without declaring
one compulsory.

\subsection{The Yang--Mills discussion does not solve Yang--Mills}

Correct.  Its relevance is diagnostic.  It identifies where locality, gauge
invariance, positivity, normalization, physical scaling, completeness, and
spectrum enter; it proves the last abstract transport step; and it locates the
first open physical estimate.  A methodology paper should not convert that
clarity into a physical theorem.  Its success is measured by whether future
proofs state and close the interfaces, not by suggesting that they are already
closed.
